\section{Introduction}

Vision-Language-Action (VLA) models unify visual perception, language understanding, and continuous action generation in a single architecture. Models such as $\pi_0$~\cite{black2024pi0}, OpenVLA~\cite{kim2024openvla}, and RT-2~\cite{brohan2023rt2} can follow diverse natural language instructions across manipulation tasks. As these models move toward real-world deployment, a critical failure mode has emerged: frequent collisions with non-target objects in cluttered environments. The demonstration data used for training typically features unobstructed workspaces, and the resulting policies acquire no explicit incentive to maintain clearance from nearby surfaces. Collisions cause task failure, risk damage to the robot and its surroundings, and represent a central barrier to safe deployment. The SafeLIBERO benchmark~\cite{hu2025vlsa}, comprising 32 collision scenarios across 1,600 evaluation episodes, was introduced to systematically quantify this problem and confirms that even high-performing VLAs exhibit substantial collision rates when non-target objects surround the manipulation target.

Several recent methods target collision avoidance in VLAs. AEGIS~\cite{hu2025vlsa} applies post-hoc Control Barrier Function (CBF) corrections to already-generated trajectories and reports an average Collision Avoidance Rate (CAR) of 77.85\% across SafeLIBERO scenarios without additional training. SafeVLA~\cite{zhang2025safevla} formulates safety as a constrained Markov Decision Process and reduces safety violations by 83.58\% through reward shaping during reinforcement learning. CoFreeVLA~\cite{cofreeevla2026} trains a risk estimator for dual-arm settings that halts execution when predicted collision risk exceeds a threshold. In our own experiments, a Pi0 model fine-tuned with an auxiliary Signed Distance Field (SDF) loss (referred to as A1-LoRA) achieves 91.5\% CAR on SafeLIBERO. These methods share a common evaluation pattern: they report high CAR as evidence of effective collision avoidance. The numbers appear compelling, and the implicit assumption across this emerging literature is that high CAR reflects genuine obstacle avoidance behavior.

This assumption is wrong. We find that the high CAR values reported by these methods are systematically contaminated by \emph{action collapse}, a phenomenon in which the VLA ceases to produce meaningful motion. When we examine the end-effector displacement of the 732 episodes that A1-LoRA classifies as ``safe,'' 99.6\% exhibit a maximum displacement below 0.1\,mm. The robot did not avoid obstacles; it avoided motion altogether. This contamination is not specific to A1-LoRA or to SDF-based approaches. The LoRA baseline (no collision avoidance objective) achieves 82\% CAR with 99.4\% of its safe episodes below the displacement threshold. Pi0.5, a more capable model from the same family, exhibits the same pattern: 98.2\% of its safe episodes are fake safe. We acknowledge that action collapse itself is a well-documented failure mode of undertrained or distribution-shifted models, and its existence alone would not warrant a dedicated analysis. What is non-trivial is the \emph{consequence}: action collapse silently inflates the primary safety metric used to evaluate collision avoidance methods. Downstream researchers who compare methods based on reported CAR values will draw incorrect conclusions about relative safety. Models that collapse more severely will paradoxically appear safer, and genuinely unsafe models may be deployed on the basis of contaminated evaluations. We propose the displacement-filtered CAR (CAR$_\text{f}$), which excludes episodes with sub-threshold displacement from both numerator and denominator, as a corrective metric that separates genuine avoidance from inaction.

The fake safety problem connects to two deeper structural issues. First, task capability and collision avoidance are anti-correlated under current evaluation practice. Pi0.5 achieves a TSR of 74.7\,$\pm$\,1.5\% but a CAR of only 26.8\,$\pm$\,1.6\%, while Pi0-based methods report CAR of 82--91.5\% alongside TSR $\leq$ 0.5\%. The explanation is mechanical: a model that moves the robot through the workspace encounters obstacles; a model that does not move encounters none. CAR is therefore informative only when conditioned on sufficient task capability, and the TSR-CAR plane reveals that no current method occupies the upper-right quadrant where both metrics are high. Second, end-effector-only SDF refinement, the most direct geometric approach to collision avoidance, has a structural defect. We construct oracle experiments that replace the learned SDF network with ground-truth signed distance fields computed from simulation geometry. All three oracle variants achieve \emph{lower} CAR than unrefined baselines (best oracle: 76.9\% vs.\ LoRA baseline: 82.0\%). The SDF gradient pushes the end-effector away from obstacles, but the resulting joint-space displacement rotates intermediate arm links toward or through obstacle volumes. Perfect distance information at the end-effector cannot compensate for the absence of full-body collision awareness. A preliminary result provides a constructive counterpoint: a Pi0.5 model augmented with training-time SDF supervision (10K-step checkpoint) achieves CAR of 56--58\% with TSR of 47--49\% and a frozen-action rate of 0\%, placing it in the previously empty central region of the TSR-CAR plane which suggests that genuine collision avoidance is achievable when the base model retains sufficient capability.

Our contributions are as follows:
\begin{itemize}
    \item We identify the fake safety problem in VLA collision avoidance evaluation: action collapse systematically inflates Collision Avoidance Rate, and we find that 98--99\% of safe episodes across multiple methods and model scales are attributable to inaction rather than obstacle avoidance. We propose displacement-filtered CAR (CAR$_\text{f}$) as a corrective metric (Section~\ref{sec:fake_safety}).
    \item We demonstrate that model capability and reported collision avoidance are anti-correlated under standard CAR evaluation. Models that execute tasks collide more because they move, while action-collapsed models receive inflated safety scores, rendering cross-method comparisons on raw CAR misleading (Section~\ref{sec:tsr_car}).
    \item Through ground-truth SDF oracle experiments, we establish that end-effector-only SDF refinement is a structural limitation: even perfect distance information cannot prevent full-body collisions, as gradients that repel the end-effector simultaneously drive arm links into obstacle volumes (Section~\ref{sec:ee_only}).
    \item[\textbf{[C4]}] Preliminary evidence from a Pi0.5 model with training-time SDF supervision shows 0\% frozen-action rate alongside CAR of 56--58\% and TSR of 47--49\%, suggesting that SDF auxiliary losses can produce genuine collision avoidance when applied to a VLA with sufficient baseline capability (Section~\ref{sec:a1_pi05}; full results pending).
\end{itemize}
